% ID: SEL-D0061-OA-E02
% ESTADO: validada. S061.
\textbf{a)} $f'(x)=4x-x^2=x(4-x)$. Es negativa en $(-\infty,0)$ y $(4,\infty)$ y positiva en $(0,4)$. La función decrece en los dos primeros intervalos de signo negativo y crece en $(0,4)$.

\textbf{b)} El cambio de signo da un mínimo relativo en $(0,0)$ y un máximo relativo en $(4,32/3)$.

\textbf{c)} La pendiente vale $4$ cuando $x(4-x)=4$, es decir, $(x-2)^2=0$. El punto es $\boxed{(2,16/3)}$.
